\documentclass[notheorems]{beamer}

\usepackage[T2A]{fontenc}
\usepackage[utf8]{inputenc}
\usepackage[english,russian]{babel}

\usetheme{Madrid}
\usecolortheme{default}
\usepackage{amssymb}
\usepackage[perpage]{footmisc}
\usepackage{amsmath}
\usepackage{graphicx}
\usepackage{amsmath,amssymb}

\newcommand{\curly}{\mathrel{\leadsto}}

\let\lemma\relax
\let\definition\relax
\let\problem\relax
\let\theorem\relax

\usepackage{amsthm}
\usepackage{comment}
\usepackage{amsmath}
\usepackage{yfonts}
\usepackage{tikz}
\usetikzlibrary{calc,patterns,angles,quotes}
\usetikzlibrary{shapes,positioning,intersections,quotes}
%\pagestyle{fancy}
\usepackage[cmtip,all]{xy}
\usepackage{graphicx}
\graphicspath{ {./} }

\usepackage{cite,caption}

% Работа с таблицами
\usepackage{array,tabularx,tabulary,booktabs}
\setbeamertemplate{caption}[numbered]

% Написание псевдокода
\usepackage{algorithm}
\usepackage{algpseudocode}
\captionsetup[algorithm]{labelsep=colon}
\floatname{algorithm}{Алгоритм}
\algrenewcommand\algorithmicrequire{\textbf{Вход:}}
\algrenewcommand\algorithmicensure{\textbf{Выход:}}

\setbeamertemplate{frametitle}[default][center]
%\usepackage[backend=biber,style=alphabetic,sorting=ynt]{biblatex}

%\addbibresource{refs.bib}

\title[Защита диплома ДМ]{Постквантовая криптография на решётках}
\author{Михаил Савватеев}
\institute[]{Научный руководитель: Алексей Фролов}
\date[22 июня 2026]{Защита диплома\\Кафедра дискретной математики МФТИ\\Долгопрудный, 22 июня 2026 г.}
\logo{}

%\AtBeginBibliography{\footnotesize}
\begin{document}

\include{include/abbreviations.tex}

\frame{\titlepage}

\begin{frame}{Виды криптографии}
\begin{itemize}
	\item<1-> Классические алгоритмы ($\mathsf{RSA}$, Диффи-Хэллман, Эль-Гамаль)
	
	Сложность факторизации и дискретного логарифмирования
	
	\smallskip
	\color{red}{Уязвимы к атаке на квантовом компьютере (алгоритмы Шора)}
	
	\medskip
	\item<2-> Изогении суперсингулярных эллиптических кривых
	
	Решение системы полиномиальных уравнений в $\F_q$
	
	\smallskip
	\color{red}{Ненадёжны (взлом $\mathsf{SIDH}$ в $2022$, $\mathsf{GeMSS}$ в $2021$, $\mathsf{Rainbow}$ в $2022$)}
	
	\medskip
	\item<3-> Хэш-функции
	
	\smallskip
	\color{cyan}{Удобны для цифровой подписи. Нет полноценного шифрования.}
	
	\medskip
	\item<4-> Коды, исправляющие ошибки
	
	\smallskip
	\color{cyan}{Изучены, но неэффективны. Многие попытки ускорить уязвимы.}
	
	\medskip
	\item<5-> Теория решёток. $\mathsf{MLWE}, \mathsf{MSIS}$ и $\mathsf{NTRU}$ problems.
	
	\smallskip
	\color{green}{Быстры, надёжны и универсальны.}
\end{itemize}
\end{frame}

\begin{frame}{LWE-problem}
Есть система зашумлённых линейных уравнений:
\begin{equation*}
	A\vec s+\vec e = \vec b \pmod{\Z_q}
\end{equation*}

\begin{itemize}
	\item $\vec s$ --- генерируется из $\chi^n$
	\item $A$ --- равномерно из $\Z_q^{m \times n}$
	\item $\vec e$ --- снова из $\chi^n$
\end{itemize}

\pause
\bigskip
Устанавливается максимально допустимая вероятность ошибки $\varepsilon$.

По известным $(A, \vec b)$ нужно c $\Pr \geq 1-\varepsilon$ восстанавливать $\vec s$.
\end{frame}

\begin{frame}{RLWE-шифрование}
\begin{minipage}{0.47\textwidth}
	\uncover<1->{
	Перейдём в $\mathcal{R}_q^n = \Z_q[x]/(x^n+1)$.
	
	Зафиксируем распределение $\chi_k$.}

	\bigskip\bigskip	
	\uncover<2->{
	Для алфавита $\Sigma_Q$ размера $Q$:
	
	\vspace{-0.2cm}
	\begin{equation*}
		\Map: \Sigma_Q^n \to \mathcal{R}_q^n
	\end{equation*}
	
	\vspace{-0.5cm}
	\begin{equation*}
		\Demap: \mathcal{R}_q^n \to \Sigma_Q^n
	\end{equation*}
	}
	
	\smallskip
	\uncover<3->{
	Взлом из решения $\mathsf{RLWE}$-problem.}
\end{minipage}
\hfill
\begin{minipage}{0.41\textwidth}
	\uncover<1->{
	Генерация ключей
	\begin{algorithmic}[1]
		\item $\bold a \leftarrow \mathcal{R}_q^n$
		\item $\bold s, \bold e \leftarrow \chi_k^n$
		\item $\bold b = \bold a\bold s+\bold e$
	\end{algorithmic}
	}
	
	\bigskip
	\uncover<2->{
	Зашифрование
	\begin{algorithmic}[1]
		\item $\bold{s'}, \bold{e'}, \bold{e''} \leftarrow \chi_k^n$
		\item $\bold u = \bold a\bold{s'}+\bold{e'}$
		\item $\bold v = \bold b\bold{s'}+\bold{e''}+\Map(m)$
	\end{algorithmic}
	}
	
	\bigskip
	\uncover<3->{
	Расшифрование
	\begin{algorithmic}[1]
		\item $\widetilde m = \Demap(\bold v-\bold u\bold s)$
	\end{algorithmic}
	}
\end{minipage}
\end{frame}

\begin{frame}{RLWE-канал}
\begin{equation*}
	\bold v-\bold u\bold s = \Map(m)+\bold z
\end{equation*}

$\mathsf{RLWE}$-зашифрование $\eqv$ передача данных по $\mathsf{RLWE}$-каналу

\pause

\vspace{0.8cm}
Использование $\mathsf{ECC}$ для компенсации зашумления.

\pause
\vspace{0.7cm}
\begin{itemize}
	\item<3-> Зависимость скорости от вероятности ошибки?
	\item<4-> Их взаимосвязь с размером алфавита?
	\item<5-> Поведение при изменении метрики в $\mathsf{ECC}$?
\end{itemize}
\end{frame}

\begin{frame}{Полученные результаты}
Криптосистемы: $\LAC\,256$ и $\NewHope\,1024$.

Целевой $\DFR \leq 10^{-6}$ (допустимое $\varepsilon = 10^{-6}$).

\pause
\bigskip
Скорость оценена через границы Хэмминга и Варшамова-Гилберта.

В $\mathsf{ECC}$ использовались метрики Хэмминга и Ли.

\only<3>{
\begin{table}[h]
	\centering
	\begin{minipage}{0.48\textwidth}
		\caption{LAC, Хэмминг}
		\centering
		\begin{tabular}{|c|c|c|}
			\hline Q & $R_{GV}$ & $R_H$ \\ \hline
			2 & 0.7569 & 0.8591 \\ \hline
			3 & 0.6298 & 0.9939 \\ \hline
		\end{tabular}
	\end{minipage}
	\hfill
	\begin{minipage}{0.48\textwidth}
		\caption{LAC, Ли}
		\centering
		\begin{tabular}{|c|c|c|}
			\hline Q & $R_{GV}$ & $R_H$ \\ \hline
			2 & 0.7569 & 0.8591 \\ \hline
			3 & 0.6298 & 0.9939 \\ \hline
		\end{tabular}
	\end{minipage}
\end{table}
}

\only<4>{
\setcounter{table}{2}
\begin{table}[h]
	\centering
	\begin{minipage}{0.48\textwidth}
		\centering
		\caption{NewHope, Хэмминг}
		\begin{tabular}{|c|c|c|}
			\hline Q & $R_{GV}$ & $R_H$ \\ \hline
			8 & 2.9342 & 2.9650 \\ \hline
			14 & 3.1555 & 3.4450 \\ \hline
			34 & 0.4060 & 2.1680 \\ \hline
		\end{tabular}
	\end{minipage}
	\hfill
	\begin{minipage}{0.48\textwidth}
		\centering
		\caption{NewHope, Ли}
		\begin{tabular}{|c|c|c|}
			\hline Q & $R_{GV}$ & $R_H$ \\ \hline
			8 & 2.9448 & 2.9703 \\ \hline
			14 & 3.3518 & 3.5442 \\ \hline
			34 & 2.8495 & 3.5972 \\ \hline
		\end{tabular}
	\end{minipage}
\end{table}
}
\end{frame}

\begin{frame}[plain]
\vfill
\begin{center}
	{\Huge Спасибо за внимание!}
	
	\vspace{1cm}
	
	Контакты: \texttt{savvateev.ma@phystech.edu}
	
	\vspace{0.2cm}
	
	\texttt{https://github.com/Dart-Xeyter/Graduate\_Thesis/}
\end{center}
\end{frame}

\end{document}

